\section{New-Build Gentrification: Its Histories, Trajectories, and Critical Geographies}

\begin{description}
  \item[Authors] Davidson, M.; Lees, L.
  \item[Year] 2010 (published online 22 September 2009)
  \item[Journal] Population, Space and Place, 16, 395--411
  \item[DOI] 10.1002/psp.584
\end{description}

\subsection*{Domain Classification}

\begin{itemize}
  \item \textbf{Primary domain(s)}: Geography
  \item \textbf{Secondary domain(s)}: Urban Economics
  \item \textbf{Keywords}: new-build gentrification, exclusionary displacement, direct chain displacement, rent gap theory, phenomenology of place, socio-economic classification (SEC), state-led gentrification, urban renaissance policy, mixed communities, right to stay put
\end{itemize}

\subsection*{Research Question}

The paper asks whether new-build residential development on brownfield and pre-existing residential land constitutes a genuine form of gentrification that produces population displacement, against the claim --- advanced by Lambert and Boddy (2002) and elaborated by Boddy (2007) under the label ``reurbanisation'' --- that such development is a class-neutral, non-displacing process of ``residentialisation.'' A secondary question, pursued through three London case studies, is how displacement associated with new-build gentrification should be conceptualized and evidenced once it is understood as more than a single ``spatial moment.''

\subsection*{Theoretical Framework}

The paper is anchored in the evolving definition of gentrification proposed by Smith (1982, 1986, 1996), who argued that gentrification had moved beyond the rehabilitation of existing housing stock to become ``the leading residential edge of a much larger endeavour: the class remake of the central urban landscape'' (Smith, 1996: 39). Against this, the authors stage a direct theoretical debate with two competing accounts of post-industrial urban demographic change: (i) the ``reurbanisation''/``residentialisation'' thesis of Lambert and Boddy (2002) and Boddy (2007), which draws on Buzar \textit{et al.}'s (2007a) second-demographic-transition framework to argue that new-build development reflects population \emph{replacement} rather than displacement; and (ii) Hamnett's (1994, 2003a,b) ``professionalisation thesis,'' which reads London's occupational upgrading as a global-city labour-market effect operating through replacement (retirement, death, out-migration of the working class) rather than gentrification-induced displacement.

To evaluate these rival accounts, the authors deploy Marcuse's (1985, 1986) four-part typology of displacement --- direct last-resident displacement, direct chain displacement, exclusionary displacement, and displacement pressure --- and extend it with a humanist/phenomenological geography of place drawn from Tuan (1977) and Relph (1977), who distinguish ``space'' from ``place'' and treat home and neighbourhood as ontologically constitutive of identity. This phenomenological move extends Hartman's (1984) ``right to stay put'' into what Davidson (2009b) calls ``the right to (make) place,'' allowing displacement to be theorized as a loss of ``structures of feeling'' and socio-spatial bonds even where no physical relocation occurs. Following Clark (2005), the authors ultimately argue for an elastic, historically and geographically informed definition of gentrification broad enough to encompass state-led and corporate new-build forms.

\textbf{Relationship to Davidson and Lees (2005)}: the article explicitly builds on and revises the authors' own earlier work. As stated in the text, Davidson and Lees (2005) had argued, using cases of new-build gentrification along London's Thames, that direct displacement could not occur on brownfield sites but that indirect displacement --- ``exclusionary displacement'' or ``price shadowing'' --- was likely. Boddy (2007) reviewed and rejected this evidence as unpersuasive. The present paper reports that ``Davidson (2008, 2009a), however, has elaborated on our earlier findings, providing detailed empirical evidence to further substantiate our claims,'' and it extends the 2005 argument in three ways: (1) it introduces cases (the Aylesbury Estate, the Pepys Estate/Aragon Tower) where redevelopment occurs on \emph{pre-existing residential} land rather than brownfield sites, so that direct displacement --- not only indirect displacement --- is documented; (2) it substantiates the 2005 indirect-displacement claim with new primary qualitative evidence (41 interviews) from three additional Thames-side wards; and (3) it develops the phenomenological theorization of displacement (via Davidson, 2009b) that was not part of the 2005 argument as described in this text.

\subsection*{Data}

\begin{itemize}
  \item \textbf{Source(s)}: (1) Documentary/secondary case material on the Aylesbury Estate (Southwark) redevelopment and the sale and conversion of Aragon Tower on the Pepys Estate (Lewisham) to Berkeley Homes, drawn from council planning documents (\textit{The Aylesbury Estate: Revised Strategy}, 2007), press reporting, and prior work by Davidson (2007, 2008, 2009a,b); (2) 41 original in-depth, semi-structured interviews with working-class residents conducted by the authors in 2004--2005; (3) 2001 UK Census data on qualifications and socio-economic classification (SEC) for London, reported at ward scale.
  \item \textbf{Geographic scope}: London, UK, with three Thames-riverside case-study wards --- Brentford, Wandsworth, and Thamesmead --- plus the Aylesbury Estate (Southwark) and the Pepys Estate (Lewisham/Deptford). Earlier comparative material (not original data) references Vancouver's Fairview Slopes and Downtown South.
  \item \textbf{Spatial unit}: Ward-level census analysis; neighbourhood/estate-level case studies; individual-level interview data.
  \item \textbf{Time period}: Interviews conducted 2004--2005; documentary and case evidence extends to the 2007 Aylesbury Revised Strategy and a 2009 postscript referencing the onset of the financial-crisis-era housing downturn.
  \item \textbf{Sample size}: N = 41 interviews across three wards; 2001 UK Census cross-tabulations of London's approximately 5.3 million adult population aged 16--74 by qualification level and socio-economic classification (Tables 1--2 in the article); approximately 2,100 new-build housing units identified across the three riverside wards.
  \item \textbf{Key variables}: Qualitative --- residents' narratives of neighbourhood change, loss of retail and community services, perceived affordability pressure, and sense of place/belonging. Quantitative/descriptive --- ward population growth (the three riverside wards grew by an average of 60\% following new-build development, with over 80\% of the added population in SEC categories 1--4, i.e.\ high-income earners); London-wide qualification levels and SEC distributions (2001 Census); a ward-level ratio of ``middle-class'' to ``working-class'' socio-economic classifications used to map the continued presence of working-class London (Figure 3).
  \item \textbf{Data limitations}: The authors explicitly acknowledge, following Slater (2009), that ``there is no statistical data available that quantifies displacement in a convincing way,'' and that available UK datasets are ``ill-suited to an `analysis of the full social complexity of individual and family circumstances''' (Newman and Wyly, 2006: 42, as cited). Aggregated census data can demonstrate population replacement but can neither prove nor disprove displacement. The interview sample necessarily captures only working-class residents who \emph{remained} in the study wards, not those already displaced.
\end{itemize}

\subsection*{Methodology}

\begin{itemize}
  \item \textbf{Empirical strategy}: This is a qualitative and conceptual critical-geography study, not a quantitative causal-inference design. It combines (1) a conceptual/theoretical literature review and critique that adjudicates between competing framings of neighbourhood change (gentrification vs.\ reurbanisation vs.\ replacement); (2) two illustrative documentary case studies of state-led estate redevelopment (Aylesbury Estate; Pepys Estate/Aragon Tower); and (3) primary qualitative fieldwork consisting of 41 semi-structured, in-depth interviews with working-class residents in three London wards, interpreted thematically through Marcuse's (1985) displacement typology and a phenomenological (Tuan/Relph) reading of place.
  \item \textbf{Identification}: There is no econometric identification strategy. Causal attribution is interpretive: the authors triangulate documentary evidence of tenure change, population turnover, and socio-economic composition with residents' first-person accounts to argue that observed changes in retail provision, affordability, and social ties are attributable to new-build gentrification rather than to unrelated processes.
  \item \textbf{Spatial treatment}: Spatial variation is handled descriptively rather than econometrically: a ward-level choropleth map (Figure 3) of the ratio of professional to working-class socio-economic classifications is used to argue against London-wide, aggregate readings of demographic change (Hamnett's professionalisation thesis) by revealing persistent working-class wards at finer spatial grain. No spatial weights matrix, spatial autoregressive model, or formal test for spatial autocorrelation is used; the modifiable-areal-unit problem is implicitly invoked (the authors argue that "reurbanisation" narratives conflate scales, treating metropolitan-scale replacement as evidence against neighbourhood-scale displacement) but not formally addressed.
  \item \textbf{Robustness checks}: In lieu of statistical robustness exercises, the authors triangulate across multiple, contextually distinct cases (a state-led compulsory estate demolition at Aylesbury; a council-to-private tower conversion at Pepys/Aragon Tower; and three riverside wards at different stages of gentrification maturity --- Brentford, Wandsworth, Thamesmead) to show that indirect and place-based displacement recurs across differing institutional and market contexts, while also re-examining Butler \textit{et al.}'s (2008) census-based evidence for middle-class growth and showing that two-thirds of the claimed growth falls within SEC sub-categories (5.1, 5.2) that include many non-middle-class occupations (counter clerks, cashiers, sales assistants, petrol-pump attendants).
  \item \textbf{Methodological innovation}: The paper's main methodological contribution is conceptual: it operationalizes ``indirect displacement'' and ``displacement pressure'' qualitatively by pairing Marcuse's (1985) typology with a phenomenological framework of place (Tuan, 1977; Relph, 1977), proposing that the presence/absence of physical relocation is an inadequate empirical test for displacement.
\end{itemize}

\subsection*{Key Findings}

\begin{enumerate}
  \item \textbf{New-build gentrification produces documented displacement, contra the reurbanisation thesis}: The authors' case evidence, taken together, is presented as repudiating Boddy's and Butler's claim that new-build/city-centre development is decoupled from working-class displacement.
  \item \textbf{Direct displacement at the Aylesbury Estate}: Under the 2007 Revised Strategy, redevelopment is projected to demolish the majority of a 2,700-unit estate and rebuild 3,200 private new-build homes and 2,000 social-rented new-build homes (meeting a 40\% social-housing planning requirement). The authors estimate this will directly displace approximately 20\% of existing tenants, price out or displace existing leaseholders who purchased under ``right to buy,'' and relocate the remaining 1,850--2,050 council tenants into a new, unfamiliar mixed-income social context.
  \item \textbf{Hybrid gentrification at Aragon Tower/Pepys Estate}: A council tower block sold in 2002 to developer Berkeley Homes for over \pounds10 million was re-clad, given five additional floors (including 14 penthouse units), gated, and renamed the ``Z apartments''; low-income council tenants were moved out and middle-income owner-occupiers moved in, illustrating a ``gentrification hybrid'' that combines redevelopment and partial retention of the existing structure.
  \item \textbf{Large-scale, high-income population influx along the Thames without a corresponding statistical signature of displacement}: Across Brentford, Wandsworth, and Thamesmead, approximately 2,100 new-build units raised ward populations by an average of 60\%, with over 80\% of the new population in high-income socio-economic classifications (SEC 1--4); yet population and census data alone do not reveal the extent of working-class displacement, which the authors argue operates in ``more silent registers.''
  \item \textbf{Indirect and place-based displacement documented via interviews}: Across the three wards, working-class interviewees described commercial/service displacement (loss of affordable local retail and amenities), exclusionary and direct-chain displacement (rising unaffordability foreclosing future occupancy, e.g.\ concern that ``new apartments would drive up prices in one of the last affordable places''), and a phenomenological loss of ``sense of place'' and socio-spatial bonds even when they had not been physically relocated. This pattern varied by context: gradual disconnection in Brentford; advanced bereavement and dislocation in the longer-gentrified Wandsworth (``this place is unrecognisable, just a few remnants like us around''); and gentrification acting as a ``disabler'' of community place-making amid pre-existing deprivation in Thamesmead.
  \item \textbf{Critique of Hamnett's professionalisation/replacement thesis}: A ward-scale re-examination of 2001 Census data (Figure 3) shows persistent and spatially concentrated working-class populations across many London wards, challenging city-wide narratives of straightforward middle-class replacement; the authors further show that two-thirds of the middle-class growth reported by Butler \textit{et al.} (2008) for 1991--2001 falls in SEC groups (5.1, 5.2) whose constituent occupations are not uniformly middle class.
\end{enumerate}

No formal statistical coefficients, effect sizes, or significance tests are reported; findings are qualitative and interpretive, supported by descriptive percentages (e.g.\ the 60\% average ward population increase, over 80\% of new residents in SEC 1--4, and the estimated 20\% direct-displacement share at Aylesbury) rather than inferential statistics.

\subsection*{Contributions}

\begin{enumerate}
  \item \textbf{Empirical contribution}: New primary qualitative evidence (41 interviews) documenting indirect and place-based displacement processes across three London new-build gentrification sites, together with updated documentary evidence on state-led estate redevelopment (Aylesbury, Pepys/Aragon Tower) that had not been part of the authors' 2005 Thames riverside argument.
  \item \textbf{Theoretical contribution}: A refinement of the concept of displacement that synthesizes Marcuse's (1985) four-part typology with humanist/phenomenological geography (Tuan, Relph), reframing displacement as a loss of place and ``structures of feeling'' rather than solely a physical, spatial moment; and an argument, following Clark (2005), for a broad, historically informed definition of gentrification inclusive of state-led and corporate new-build forms.
  \item \textbf{Policy contribution}: A critique of UK New Labour ``urban renaissance'' and ``mixed communities''/``positive gentrification'' policy (echoing the US HOPE VI program) as a mechanism that produces displacement under the rhetoric of social mixing, with the Aylesbury Estate cited as a paradigm case.
\end{enumerate}

No distinct methodological contribution in the sense of a novel identification strategy is claimed; the paper's innovation is conceptual/qualitative rather than statistical.

\subsection*{Limitations and Critical Assessment}

\begin{enumerate}
  \item \textbf{Identification threats}: The qualitative case-comparison design cannot rule out confounding from city-wide or macroeconomic forces unrelated to new-build gentrification specifically --- for example, London-wide retail restructuring, cost-of-living increases, or an aging housing stock could independently generate the same resident narratives of loss and unaffordability that the authors attribute to gentrification. The 41 interviews are drawn from residents who \emph{remained} in the study wards, introducing a survivorship bias: those already displaced (arguably the population of greatest interest) are structurally absent from the sample, which limits the paper's ability to speak to the magnitude of direct or chain displacement it also claims is occurring.
  \item \textbf{External validity}: The case studies are purposively selected within London, an exceptional global-city housing market with a distinctive tenure structure (large council-housing stock, ``right to buy'' policy, New Labour urban renaissance planning framework). The phenomenological account of place-loss, while theoretically portable, is evidenced only in this institutional and cultural context; generalization to other national housing systems (e.g., weaker social-housing sectors, different eviction-protection regimes, or non-Anglophone contexts) is not established and is not tested by the authors.
  \item \textbf{Omitted mechanisms}: The paper does not trace where displaced tenants and leaseholders actually relocate, nor does it assess the welfare consequences (housing quality, commute, social-network disruption) of relocation destinations, despite adopting a place-based welfare framework that would seem to require this. It also does not examine the financial mechanisms underlying redevelopment decisions --- for instance, land-value uplift capture, developer viability assessments, or the specific terms of the Berkeley Homes/Aragon Tower transaction --- even though the paper's own narrative foregrounds ``corporate capital'' and state asset transfer as drivers of new-build gentrification; a geographic-finance-style analysis of deal economics is absent.
  \item \textbf{Data constraints}: The authors correctly note that no UK administrative or census data source can quantify displacement directly. However, the paper does not attempt to approximate displacement magnitude using available proxies (e.g., electoral-roll churn, housing-benefit or council-tax records, or the New York City Housing and Vacancy Survey-style approach that Newman and Wyly (2006) used to bound displacement rates at 6.2--9.9\% of renter households, a study the authors cite only in a footnote). Linking the 41 interviews to household-level administrative panel data on relocation would substantially strengthen the causal narrative.
  \item \textbf{Equilibrium considerations}: The analysis operates entirely at the neighbourhood/estate scale and does not consider general-equilibrium housing-market responses. Notably, the paper does not address whether the substantial new housing supply it documents (2,100 units along the Thames; 3,200 private units at Aylesbury) could, in aggregate, relieve city-wide affordability pressure even while displacing specific incumbent communities locally --- an omission that matters for policy interpretation, since a pure neighbourhood-scale critique of new-build development as ``a displacement machine'' (a framing the paper itself attributes critically to Hamnett's logic) could equally be levelled, on the paper's own terms, against any housing-supply expansion in a supply-constrained city.
\end{enumerate}

\subsection*{Cross-Disciplinary Bridges}

\begin{itemize}
  \item \textbf{Connection to Urban Economics}: The paper's evidence of large, income-selective in-migration following new-build supply additions (over 80\% of new residents in SEC 1--4) speaks directly to urban-economic debates on housing supply elasticity (Glaeser and Gyourko) and filtering theory. An urban-economics reading would ask whether the exclusionary displacement the authors document reflects amenity capitalization in the Rosen (1974) hedonic sense (new-build attributes bid up local prices) or a supply-side response that, in general equilibrium, could offset city-wide affordability pressure even as it disrupts specific incumbent communities.
  \item \textbf{Connection to Geographic Finance}: The state-led transfer of public housing assets to corporate developers (Berkeley Homes at Aragon Tower; the \pounds2.4 billion Aylesbury redevelopment) is a clear instance of the financialization of housing (Aalbers, 2016) --- public land and buildings recommodified as vehicles for private capital accumulation. A geographic-finance lens would examine the underlying land-value capture mechanics, developer return targets, and financing structure of these deals, none of which the present paper analyzes.
  \item \textbf{Connection to Regional Science}: The paper's ward-level remapping of London's socio-economic geography (Figure 3), used to contest a metropolitan-scale ``replacement'' narrative, is a direct illustration of the modifiable areal unit problem (MAUP) and of scale-dependent inference familiar to regional science, even though the authors do not use that terminology.
  \item \textbf{Suggested related works}:
    \begin{itemize}
      \item Freeman, L. (2005) --- \textit{Displacement or succession? Residential mobility in gentrifying neighborhoods}, \textit{Urban Affairs Review}: offers a competing quantitative/survey-based measurement of displacement that this paper directly critiques as understating the phenomenon; a useful methodological counterpoint.
      \item Aalbers, M. (2016) --- \textit{The Financialization of Housing}: provides the theoretical apparatus for a geographic-finance interpretation of the state-to-corporate asset transfers documented at Aragon Tower and the Aylesbury Estate.
      \item Saiz, A. (2010) --- \textit{The geographic determinants of housing supply}, \textit{Quarterly Journal of Economics}: an urban-economics/regional-science counterpoint on housing supply constraints, relevant to testing whether the new-build supply documented here is best understood as relieving or intensifying London's affordability crisis.
      \item Chetty, R., Hendren, N.\ and Katz, L.F.\ (2016) --- \textit{The effects of exposure to better neighborhoods on children: New evidence from the Moving to Opportunity experiment}, \textit{American Economic Review}: a contrasting, experimentally identified treatment of neighbourhood change and residential mobility, useful for benchmarking the evidentiary standard of this paper's phenomenological approach.
    \end{itemize}
\end{itemize}

\subsection*{Quotable Passages}

\begin{quote}
"Gentrification is no longer about a narrow and quixotic oddity in the housing market but has become the leading residential edge of a much larger endeavour: the class remake of the central urban landscape" (p.~396, quoting Smith, 1996: 39).
\end{quote}

\begin{quote}
"[D]isplacement affects many more than those actually displaced at any given moment" (p.~408, quoting Marcuse, 1986: 157).
\end{quote}

\begin{quote}
"We argue that a purely spatial account of displacement is inadequate. As such, we understand displacement as operating uniquely across neighbourhoods, according to the particular contexts and positionalities" (p.~408).
\end{quote}
